Parkinson's disease (PD) is a chronic, progressive neurodegenerative disorder primarily characterized by the substantial loss of dopaminergic neurons in the substantia nigra
\citep{Nguyen:2021, Fujikawa:2023, Meng:2022}.
This neurological decline results in a set of cardinal motor manifestations, including bradykinesia (slowness of movement), rigidity (stiffness), postural instability, and a characteristic unilateral rest tremor
\citep{Nguyen:2021, Botzel:2014}. 

Pharmacological therapy remains the first-line approach for tremor management, primarily relying on dopamine replacement through Levodopa or dopamine agonists \citep{Botzel:2014, Lora-Millan:2021}. However, long-term pharmacotherapy often leads to reduced effectiveness and may induce side effects including dyskinesia and motor fluctuations \citep{Lora-Millan:2021, Meng:2022}. In addition, a significant portion of patients does not respond adequately to medication \citep{Fujikawa:2023}. Surgical interventions such as Deep Brain Stimulation are considered in severe or refractory cases \citep{Botzel:2014, Fujikawa:2023}, although their invasive nature and associated risks limit their applicability. Complementary strategies, including occupational therapy and assistive tools, are also adopted to mitigate functional impairment \citep{Mo:2021, Bhidayasiri:2022}.

Orthotic devices are a less invasive alternative and can be classified as passive or active. Passive orthoses rely on mechanical elements or smart materials to attenuate tremor by dissipating vibratory energy at specific frequencies \citep{Gebai:2018, Moura:2025}, but are sensitive to variations in tremor frequency and require further validation in real-world conditions \citep{Heide:2014, Bhidayasiri:2022}. Active orthoses, in contrast, employ closed-loop control strategies that estimate tremor characteristics in real time and apply compensatory actions through mechanical loading or electrical stimulation \citep{Lora-Millan:2021, Meng:2022}.

Active control approaches introduce additional challenges due to their reliance on sensing, actuation, and real-time computation, which may reduce user comfort and cause fatigue \citep{Lora-Millan:2021, Martinez-Nunez:2024}. In particular, peripheral or functional electrical stimulation techniques may also induce discomfort and skin irritation or fatigue \citep{Taheri:2015, Saleh:2022}. Despite this, advances in hardware miniaturization and signal processing ensure their relevance, especially in wrist tremor suppression.

Several control strategies for wrist tremor suppression have been proposed in the literature. \cite{Gallego:2013} design a neuroprosthesis based on Proportional Integrative (PI) control for a single degree of freedom (DoF) system of the human wrist, combining a Critically Damped g-h Filter (CDF), a Weighted Frequency Fourier Linear Combiner (WFLC), and a Kalman filter for tremor estimation. \cite{Copur:2019} propose a zero-phase high-pass filter for tremor estimation alongside a repetitive control strategy for electrical stimulation-based co-contraction in a single DoF setup.

In parallel, \cite{Lekshmi:2019} introduce a Proportional Integrative Derivative (PID) controller applied to a four DoF wrist model with gains optimized via genetic algorithm. \cite{Zamanian:2019} develop an adaptive notch filter supported by an Adaptive Frequency Estimator (AFE) to track tremor frequency in the velocity signal. More recently, \cite{Qi:2024} propose an Active Disturbance Rejection Control (ADRC) strategy using an Enhanced Bandlimited Multiple Fourier Linear Combiner (EBMFLC) for voluntary motion estimation.

Most approaches focus on single-joint models, where excellent results have been achieved, but do not account for dynamic coupling between joints, where shoulder and elbow PD tremors significantly affect wrist motion. To address this issue, a 3-DoF linear upper limb model proposed in \citep{Gebai:2017, Gebai:2018} is adopted in this work as a test platform.

In this scenario, disturbance rejection-based control methods are expected to offer advantages due to their robustness to model uncertainties and external perturbations. This work investigates an error-based active disturbance rejection control for wrist tremor suppression under coupled dynamics. Two variations are considered, differing in how voluntary motion is estimated: one based on EBMFLC (following after \cite{Qi:2024}) and another on Zero-Phase Low-Pass (ZPLP) filtering. 
These are compared against AFE,  
the PI proposed by \cite{Gallego:2013}, and two PID controllers, one tuned with Internal Model Control (IMC) and the other with offline Differential Evolution (DE) optimization. 
Robustness is evaluated through Monte Carlo simulations of physical parameters.

The remainder of this paper is organized as follows: Section \ref{sec:theory} presents the upper limb dynamic model and the control framework. Section \ref{sec:methodology} describes the simulation setup. Section \ref{sec:results} discusses the obtained results. Section \ref{sec:conclusion} concludes the paper and outlines future work.